\chapter{Conclusi\'on general}\label{chapter:conclusion}

\section{Resumen de aportes}

MiGanado propone una plataforma integral para establecimientos bovinos de cr\'ia para la venta que combina registro operativo, organizaci\'on multiestablecimiento, gesti\'on sanitaria y reproductiva, calendario de alertas, indicadores, asistencia inteligente, detecci\'on de anomal\'ias y proyecciones econ\'omicas. Su aporte principal consiste en integrar funcionalidades que habitualmente aparecen separadas o con alcance parcial, orient\'andolas a un contexto productivo espec\'ifico de la provincia de Buenos Aires.

Desde el punto de vista acad\'emico y profesional, el proyecto permite aplicar conocimientos de ingenier\'ia de software, dise\~no de sistemas, modelado de datos, experiencia de usuario, inteligencia artificial y validaci\'on de prototipos en un dominio real con impacto productivo.

\section{Trabajo futuro}

Como trabajo futuro se prev\'e ampliar la cobertura geogr\'afica del sistema a otras regiones de Argentina, incorporar integraciones con hardware externo de identificaci\'on o lectura animal, profundizar los modelos predictivos con datasets curados, ampliar el m\'odulo econ\'omico, validar el chatbot con usuarios del dominio y realizar pruebas piloto en establecimientos reales.

Tambi\'en quedar\'a pendiente completar las secciones de investigaci\'on de usuarios, anexar encuestas y entrevistas, incorporar diagramas definitivos, documentar tecnolog\'ias seleccionadas, definir datasets y consolidar la bibliograf\'ia con fuentes acad\'emicas y t\'ecnicas adicionales.
